\documentclass[12pt,a4paper]{report}
\usepackage[utf8]{inputenc}
\usepackage[T1]{fontenc}
\usepackage[french]{babel}
\usepackage{geometry}
\geometry{margin=2.5cm}
\usepackage{hyperref}
\usepackage{listings}
\usepackage{xcolor}
\usepackage{graphicx}
\usepackage{booktabs}
\usepackage{longtable}
\usepackage{array}
\usepackage{fancyhdr}
\usepackage{titlesec}

% Screenshot macro with title and image
\newcommand{\screenshot}[2]{%
  \begin{figure}[h!]
    \centering
    \includegraphics[width=0.8\textwidth]{#2}
    \caption{#1}
    \label{fig:#2}
  \end{figure}
}

\definecolor{codegreen}{rgb}{0,0.6,0}
\definecolor{codegray}{rgb}{0.5,0.5,0.5}
\definecolor{codepurple}{rgb}{0.58,0,0.82}
\definecolor{backcolour}{rgb}{0.95,0.95,0.92}

% FIX 1: Define Dockerfile as a listings language (not built-in)
\lstdefinelanguage{Dockerfile}{
  keywords={FROM, RUN, CMD, LABEL, EXPOSE, ENV, ADD, COPY, ENTRYPOINT,
            VOLUME, USER, WORKDIR, ARG, ONBUILD, STOPSIGNAL, HEALTHCHECK, SHELL},
  keywordstyle=\color{magenta},
  sensitive=true,
  comment=[l]{\#},
  commentstyle=\color{codegreen},
  stringstyle=\color{codepurple},
  morestring=[b]",
  morestring=[b]',
}

\lstdefinestyle{mystyle}{
    backgroundcolor=\color{backcolour},
    commentstyle=\color{codegreen},
    keywordstyle=\color{magenta},
    numberstyle=\tiny\color{codegray},
    stringstyle=\color{codepurple},
    basicstyle=\ttfamily\footnotesize,
    breaklines=true,
    numbers=left,
    numbersep=5pt,
    showspaces=false,
    showstringspaces=false,
    tabsize=2
}
\lstset{style=mystyle}

\pagestyle{fancy}
\fancyhf{}
\fancyhead[L]{Déploiement d'une Application Full-Stack}
\fancyhead[R]{\leftmark}
\fancyfoot[C]{\thepage}

\titleformat{\chapter}[display]
  {\normalfont\huge\bfseries}{\chaptertitlename\ \thechapter}{20pt}{\Huge}
\titleformat{\section}
  {\normalfont\Large\bfseries}{\thesection}{1em}{}
\titleformat{\subsection}
  {\normalfont\large\bfseries}{\thesubsection}{1em}{}

\begin{document}

% ==============================
% PAGE DE GARDE
% ==============================
\begin{titlepage}
\centering
\vspace*{2cm}
{\Huge\bfseries Déploiement d'une Application Web Full-Stack}\\[0.5cm]
{\LARGE\bfseries Docker Compose \& Kubernetes (K3s)}\\[2cm]
{\Large Application : Simple 3-Tier Tasks Web App}\\[0.3cm]
{\large Stack : React / Express.js / MongoDB / Nginx}\\[3cm]
\rule{\textwidth}{1pt}\\[0.5cm]
{\large\itshape Rapport technique de projet DevOps - K3s}\\[0.5cm]
\rule{\textwidth}{1pt}\\[3cm]
{\large\textbf{Outils :}} Docker/Compose CLI, kubectl, K3s, VS Code, Git\\[2cm]
{\large Date : Mai 2026}
\vfill
\end{titlepage}

% FIX 2: Add table of contents (was missing)
\tableofcontents
\clearpage

% ==============================
% RÉSUMÉ
% ==============================
\chapter*{Résumé / Executive Summary}
\addcontentsline{toc}{chapter}{Résumé / Executive Summary}

Ce rapport présente le déploiement d'une application web trois-tiers (gestionnaire de tâches) avec Docker Compose pour le développement local, puis Kubernetes (K3s) pour la production sur un cluster de trois nœuds. L'application a été choisie pour sa simplicité, permettant de se concentrer sur les bonnes pratiques DevOps.

\textbf{Technologies :} React (Vite) + Express.js + MongoDB + Nginx, orchestrées par Docker Compose et K3s.

\textbf{Résultats :} Déploiement fonctionnel avec multi-stage builds, utilisateurs non-root, limites de ressources, healthchecks, probes, secrets externalisés, NetworkPolicies, Ingress et HPA.

% ==============================
% INTRODUCTION
% ==============================
\chapter{Introduction}

\section{Présentation de l'application}

L'application choisie est un \textbf{gestionnaire de tâches simple} (3-Tier Tasks Web App). Sa simplicité délibérée permet de se concentrer sur la complexité du déploiement.

\textbf{Fonctionnalités :} CRUD de tâches via interface React, API REST Express.js, persistance MongoDB.

\section{Architecture globale}

Architecture trois-tiers :
\begin{enumerate}
\item \textbf{Frontend} : React servi par Nginx (port 80/8080)
\item \textbf{Backend} : API REST Express.js (port 3000)
\item \textbf{Database} : MongoDB (port 27017)
\end{enumerate}

\section{Choix technologiques}
\begin{itemize}
\item \textbf{React + Vite} : build rapide, assets optimisés
\item \textbf{Express.js} : backend léger et performant
\item \textbf{MongoDB} : NoSQL documentaire adapté aux données flexibles
\item \textbf{Nginx} : serveur statique et reverse proxy
\item \textbf{Docker Compose} : orchestration locale simple
\item \textbf{K3s} : Kubernetes allégé pour environnements limités
\end{itemize}

% ==============================
% ENVIRONNEMENT
% ==============================
\chapter{Préparation de l'environnement}

\section{Outils installés}
\begin{itemize}
\item Docker Engine et Docker Compose CLI
\item kubectl (client Kubernetes)
\item K3s (distribution allégée)
\item VS Code (édition des manifests)
\item Git (versioning)
\end{itemize}

\section{Configuration de la machine hôte}

Machine hôte Windows 16~Go RAM, 3 VMs Debian sous VMware Workstation :

\begin{table}[h]
\centering
\begin{tabular}{|l|l|l|l|}
\hline
\textbf{VM} & \textbf{Rôle} & \textbf{IP host-only} & \textbf{IP bridged} \\
\hline
k3s-cp & Control-plane & 192.168.204.129 & 192.168.1.20 \\
k3s-w1 & Worker 1 & 192.168.204.130 & -- \\
k3s-w2 & Worker 2 & 192.168.204.131 & -- \\
\hline
\end{tabular}
\caption{Configuration des VMs}
\end{table}

\section{Installation du cluster K3s}

\subsection{Topologie réseau}
\begin{itemize}
\item \textbf{Host-only} (192.168.204.0/24) : communication interne cluster (Flannel, K3s)
\item \textbf{Bridged} (192.168.1.0/24) : accès LAN externe pour le control-plane
\end{itemize}

\subsection{Configuration control-plane}
\texttt{/etc/rancher/k3s/config.yaml} :
\begin{lstlisting}[language=yaml]
node-ip: 192.168.204.129
node-external-ip: 192.168.1.20
flannel-iface: <interface-host-only>
tls-san:
  - 192.168.204.129
  - 192.168.1.20
\end{lstlisting}

\subsection{Workers}
Les workers rejoignent via le control-plane host-only :
\begin{lstlisting}[language=yaml]
node-ip: 192.168.204.130
flannel-iface: <interface-host-only>
\end{lstlisting}

\subsection{Difficultés rencontrées}
\begin{itemize}
\item \textbf{Clonage VMs} : modification hostname, IP, /etc/machine-id, clés SSH
\item \textbf{Mauvaise interface} : configuration explicite de \texttt{node-ip} et \texttt{flannel-iface}
\item \textbf{Changement de sous-réseau} : mise à jour IP control-plane, régénération token, réinstallation agents
\end{itemize}

\section{Versioning (Git)}
Branche \texttt{main} pour le code stable, fichiers sensibles exclus via \texttt{.gitignore} et \texttt{.dockerignore}.

% ==============================
% PHASE 1 : DOCKER COMPOSE
% ==============================
\chapter{Phase 1 : Déploiement avec Docker Compose}

\section{Structure du projet}

\begin{lstlisting}[language=bash]
.
|-- react-tasks-frontend/      # React (Vite) + nginx.conf
|-- express-tasks-backend/     # Express + Dockerfile + init/
|-- docker-compose.yml
|-- .env
|-- .dockerignore
|-- screenshots/
\end{lstlisting}

\section{Best practices appliquées}

\subsection{Versions explicites des images}
Tag \texttt{mongo:7.0} utilisé au lieu de \texttt{latest} pour garantir la reproductibilité.

\subsection{Réseaux personnalisés}
\begin{itemize}
\item \textbf{frontend-net} : frontend $\leftrightarrow$ backend
\item \textbf{backend-net} : backend $\leftrightarrow$ database
\end{itemize}
L'isolation empêche le frontend d'atteindre directement la base de données.

\subsection{Volumes nommés}
\begin{itemize}
\item \textbf{mongo-data} : volume nommé pour la persistance des données
\item \textbf{Bind mount} : script d'initialisation monté en lecture seule (\texttt{:ro})
\end{itemize}

\subsection{Variables d'environnement et fichier .env}

Toutes les variables sensibles sont externalisées dans \texttt{.env} :

\begin{lstlisting}[language=bash]
DB_ROOT_USER=...
DB_ROOT_PASSWORD=...
DB_NAME=...
DB_APP_USER=...
DB_APP_PASSWORD=...
PORT=3000
VITE_API_URL=/api
\end{lstlisting}

Le fichier \texttt{.env} est exclu de Git via \texttt{.gitignore}.

\subsection{Healthchecks et dépendances}
\begin{itemize}
\item Database : mongosh avec \texttt{db.runCommand(\{ping: 1\})}
\item Backend : HTTP /health (via Dockerfile)
\item Frontend : HTTP / (via Dockerfile)
\item \texttt{depends\_on} avec \texttt{condition: service\_healthy} pour l'ordre de démarrage
\end{itemize}

\subsection{Limites de ressources}

\begin{table}[h]
\centering
\begin{tabular}{|l|c|c|c|}
\hline
\textbf{Service} & \textbf{CPU} & \textbf{Mémoire} & \textbf{Restart} \\
\hline
Database & 0.5 & 512m & unless-stopped \\
Backend & 0.5 & 256m & unless-stopped \\
Frontend & 0.25 & 128m & unless-stopped \\
\hline
\end{tabular}
\caption{Limites de ressources Docker Compose}
\end{table}

\subsection{Commandes utilisées}

\begin{lstlisting}[language=bash]
# Build et demarrage
docker compose up -d --build

# Verification des services
docker compose ps
docker compose logs -f backend

# Statistiques runtime
docker stats
\end{lstlisting}

% ==============================
% PHASE 2 : KUBERNETES K3S
% ==============================
\chapter{Phase 2 : Déploiement sur Cluster Kubernetes (K3s)}

\section{Architecture du cluster K3s}

Le cluster est composé de 3 nœuds Debian sous VMware Workstation :

\begin{table}[h]
\centering
\begin{tabular}{|l|l|l|l|}
\hline
\textbf{Nœud} & \textbf{Rôle} & \textbf{IP interne} & \textbf{État} \\
\hline
k3s-cp & Control-plane & 192.168.204.129 & Ready \\
k3s-w1 & Worker 1 & 192.168.204.130 & Ready \\
k3s-w2 & Worker 2 & 192.168.204.131 & Ready \\
\hline
\end{tabular}
\caption{État du cluster K3s}
\end{table}

\begin{table}[h]
\centering
\begin{tabular}{|l|p{8cm}|}
\hline
\textbf{Best Practice} & \textbf{Application} \\
\hline
Namespace dédié & tasks-app \\
Labels cohérents & app.kubernetes.io/name, app.kubernetes.io/component \\
Resource requests/limits & CPU et mémoire sur tous les conteneurs \\
Liveness/Readiness probes & HTTP pour frontend/backend, exec pour MongoDB \\
RollingUpdate & maxSurge=1, maxUnavailable=0 \\
Utilisateur non-root & runAsUser: 1001 (app), 999 (MongoDB) \\
Replicas critiques & 2 pour frontend et backend \\
Secrets externalisés & db-secrets injectés via env vars \\
ConfigMaps & app-config, frontend-config, mongo-init-script \\
PVC + StorageClass & mongo-data (5Gi, app-local-path) \\
Ingress & Traefik sur tasks.local \\
HPA & backend scale 2--5 replicas selon CPU 70\% \\
NetworkPolicies & Isolation frontend/backend/database \\
\hline
\end{tabular}
\caption{Best practices appliquées sur K3s}
\end{table}

\section{Commandes d'application des manifests}

\begin{lstlisting}[language=bash]
# Application sequentielle des manifests
kubectl apply -f k3s/namespace.yaml
kubectl apply -f k3s/storageclass.yaml
kubectl apply -f k3s/secrets/secret.yaml
kubectl apply -f k3s/configmaps/configmap.yaml
kubectl apply -f k3s/deployments/
kubectl apply -f k3s/services/
kubectl apply -f k3s/ingress.yaml
kubectl apply -f k3s/backend-hpa.yaml
kubectl apply -f k3s/networkpolicy.yaml

# Verification
curl http://<node-ip>:30080  # Acces frontend via NodePort
\end{lstlisting}

\section{Vérification des déploiements}

\subsection{3 Deployments distincts (ou StatefulSet pour la DB)}

\screenshot{Replicas}{wrk_ld.jpeg}

\begin{lstlisting}[language=bash]
# Vérification des déploiements
kubectl get deployments -n tasks-app
kubectl get statefulsets -n tasks-app
kubectl get pods -n tasks-app -o wide
\end{lstlisting}

\subsection{Resource requests et limits définis sur tous les conteneurs}

\screenshot{Configuration des ressources}{res_limit.jpeg}

\begin{lstlisting}[language=bash]
# Vérification des ressources
kubectl describe pod <pod-name> -n tasks-app | grep -A 10 "Limits\|Requests"
kubectl top pods -n tasks-app
\end{lstlisting}

\subsection{Liveness Probe et Readiness Probe configurées}

\screenshot{Liveness et Readiness Probes}{rdy_liv.jpeg}

\begin{lstlisting}[language=bash]
# Vérification des probes
kubectl describe pod <pod-name> -n tasks-app | grep -A 5 "Readiness\|Liveness"
kubectl get events -n tasks-app --field-selector involvedObject.name=<pod-name>
\end{lstlisting}

\subsection{Stratégie de déploiement RollingUpdate bien paramétrée}

\screenshot{Configuration RollingUpdate}{rollout.jpeg}

\begin{lstlisting}[language=bash]
# Vérification de la stratégie de déploiement
kubectl describe deployment backend -n tasks-app | grep -A 5 "Strategy"
kubectl rollout status deployment/backend -n tasks-app
kubectl rollout history deployment/backend -n tasks-app
\end{lstlisting}

\subsection{Conteneurs exécutés en non-root (securityContext)}

\screenshot{SecurityContext - Utilisateur non-root}{sec_con.jpeg}

\begin{lstlisting}[language=bash]
# Vérification de l'utilisateur non-root
kubectl exec -it <pod-name> -n tasks-app -- whoami
kubectl exec -it <pod-name> -n tasks-app -- id
kubectl get pod <pod-name> -n tasks-app -o jsonpath='{.spec.securityContext}'
\end{lstlisting}

\section{Vérification des Services et Exposition}

\subsection{Services ClusterIP pour backend et base de données}

\screenshot{Services ClusterIP}{0.jpeg}

\begin{lstlisting}[language=bash]
# Vérification des services ClusterIP
kubectl get services -n tasks-app
\end{lstlisting}

\subsection{Service pour exposer le frontend (NodePort / LoadBalancer ou Ingress)}

\screenshot{Services, engress et configmap}{1.jpeg}

\begin{lstlisting}[language=bash]
# Vérification du service d'exposition
kubectl get service frontend -n tasks-app -o wide

kubectl get ingress -n tasks-app

# Vérification des ConfigMaps
kubectl get configmaps -n tasks-app
\end{lstlisting}

\subsection{Secrets pour tous les éléments sensibles}

\screenshot{Secret Base de données - Identifiants sécurisés}{2.jpeg}

\begin{lstlisting}[language=bash]
# Vérification des Secrets
kubectl describe secret db-secrets -n tasks-app
\end{lstlisting}

% ==============================
% PLACEHOLDERS CHAPITRES RESTANTS
% ==============================
\chapter{Tests et Validation}
\textbf{[Section à compléter par l'étudiant]}

Cette section doit présenter les tests fonctionnels et de validation effectués sur les deux environnements (Docker Compose et K3s). Elle doit inclure :

\begin{itemize}
\item Tests fonctionnels sur Docker Compose (endpoints CRUD, healthchecks, persistance)
\item Tests sur le cluster K3s (curl, navigateur, NodePort 30080, Ingress tasks.local)
\item Vérification de la haute disponibilité (arrêt d'un worker, résilience des pods)
\item Monitoring basique avec \texttt{kubectl get all}, \texttt{logs}, \texttt{describe}
\item Captures d'écran et preuves de fonctionnement
\end{itemize}

\chapter{Difficultés rencontrées et solutions apportées}
\textbf{[Section à compléter par l'étudiant]}

Cette section doit documenter les problèmes techniques rencontrés et les solutions mises en œuvre. Exemples de sujets à traiter :

\begin{itemize}
\item Problèmes de réseau lors du clonage des VMs (hostname, IP, machine-id)
\item Mauvaise sélection d'interface par K3s (node-ip, flannel-iface)
\item Rejoindre le cluster après changement de sous-réseau (token, réinstallation agents)
\item Problèmes de healthchecks et timeouts
\item Persistance des données et volumes
\end{itemize}

\chapter{Conclusion et perspectives}
\textbf{[Section à compléter par l'étudiant]}

Cette section doit inclure :

\begin{itemize}
\item Ce qui a été appris durant le projet
\item Les compétences acquises en Docker, Kubernetes et DevOps
\item Améliorations possibles : CI/CD (GitHub Actions, GitLab CI), monitoring avancé (Prometheus/Grafana), GitOps (ArgoCD), cert-manager pour TLS, multi-réplica MongoDB (replica set)
\end{itemize}

% ==============================
% ANNEXES
% ==============================
\chapter*{Annexes}
\addcontentsline{toc}{chapter}{Annexes}

\section*{Arborescence complète du projet}
\textbf{[À compléter]}

\section*{Captures d'écran importantes}
\textbf{[À insérer]}

\section*{Commandes utiles}
\textbf{[Résumé des commandes kubectl et docker les plus utilisées]}

\end{document}
